\chapter{Introducción}
Detectar un archivo malicioso es solo una parte del trabajo de un equipo de seguridad. Cuando una alerta puede llevar al aislamiento de un equipo o a la interrupción de un servicio, también interesa examinar qué elementos influyeron en la decisión. Los modelos de aprendizaje automático permiten reconocer patrones en los datos, pero su salida no siempre ofrece una justificación que pueda revisarse. Las técnicas de inteligencia artificial explicable, o XAI, buscan hacer accesible parte de ese comportamiento \parencite{Manthena2025,Liu2025GlassBox}.

En la clasificación de malware, una explicación puede mostrar la influencia de permisos, llamadas a funciones o atributos de un ejecutable. Ese resultado ayuda a inspeccionar el modelo, aunque no demuestra por sí solo qué hace el programa ni reemplaza un análisis forense. También introduce trabajo adicional: generar perturbaciones, volver a consultar el clasificador o calcular atribuciones requiere tiempo y memoria. Por eso, un sistema que predice con rapidez no necesariamente explica con la misma rapidez.

El interés de este seminario surge de esa diferencia. Se busca estudiar cuánto cuesta generar explicaciones y qué decisiones de diseño permiten reducir ese costo conservando información útil para examinar las predicciones. La pregunta tiene sentido incluso cuando el clasificador ya es eficiente, porque el comportamiento del explicador depende de su configuración, de la representación de entrada y del modelo que consulta.

La ampliación bibliográfica obliga a precisar la propuesta inicial. Kim y colaboradores emplean explicabilidad para seleccionar atributos de grafos; Ardito y colaboradores integran reducción, clasificación y auditoría mediante SHAP y LIME; y RFS-MD explora explicaciones basadas en reglas \parencite{Kim2026FCG,Ardito2026Hybrid,Alhawari2026RFSMD}. Estos antecedentes impiden sostener que la integración de eficiencia y explicabilidad carece de soluciones previas. A la vez, sus objetivos y mediciones no equivalen a una comparación controlada del costo incremental de explicar con selección, destilación y asignación adaptativa de recursos.

El proyecto propone estudiar esa combinación. Primero se seleccionarán características para reducir la entrada del modelo. Luego se evaluará si un estudiante destilado conserva suficiente capacidad predictiva frente a un profesor y frente al mismo estudiante entrenado directamente. Finalmente se asignarán distintos presupuestos de explicación, atendiendo a la incertidumbre y a criterios de auditoría. El ahorro deberá comprobarse: mostrar menos variables o dejar casos sin explicar no será suficiente para demostrar una mejora.

La investigación mantendrá un alcance experimental. Los resultados sobre conjuntos públicos permitirán comparar configuraciones bajo condiciones controladas, pero no probarán por sí solos que el sistema funciona en un SOC real o que mejora las decisiones de sus analistas. Esta distinción es pertinente porque la calidad técnica de una explicación y su utilidad para una persona son dimensiones relacionadas, pero diferentes \parencite{Cecchi2026Connecting,Tayari2026Clarity}.

El informe presenta los fundamentos y el diseño experimental. El anexo vincula las 43 fichas con la propuesta y distingue su grado de respaldo.
